% Sample Complexity Theorem (Condensed for Main Body)
% Full proof and analysis in Appendix

\subsection{Sample Complexity}
\label{subsec:sample_complexity}

\begin{theorem}[Sample Complexity for O/R Estimation]
\label{thm:sample_complexity}
For any $\epsilon > 0$ and $\delta \in (0,1)$, with $n \geq \frac{2}{\epsilon^2} \ln(4/\delta)$ i.i.d.\ trajectory samples from a policy with bounded action probabilities and non-degenerate marginal variance, we have with probability $\geq 1 - \delta$:
\begin{equation}
|\widehat{\text{OR}}_n(\pi) - \text{OR}(\pi)| \leq \epsilon
\end{equation}
\end{theorem}

\textbf{Practical implications:} For $\epsilon = 0.1$ (10\% error) at 99\% confidence, we need $n \geq 1{,}200$ trajectories—matching our experimental setup. Even $n = 30$ achieves $\epsilon \approx 0.26$, explaining why O/R predicts performance as early as episode 10. Unlike mutual information (exponential in agent count) or graph metrics (requiring global information), O/R scales linearly with precision. Full proof and caveats in Appendix~\ref{app:sample_complexity}.
